\documentclass[11pt,a4paper]{article}

\usepackage[utf8]{inputenc}
\usepackage[T1]{fontenc}
\usepackage[ngerman]{babel}
\usepackage{lmodern}
\usepackage{amsmath,amssymb,mathtools}
\usepackage{booktabs,array}
\usepackage{enumitem}
\usepackage{geometry}
\usepackage{hyperref}

\geometry{left=25mm,right=25mm,top=24mm,bottom=27mm}
\setlength{\parindent}{0pt}
\setlength{\parskip}{0.6em}
\setlist[itemize]{itemsep=0.2em,topsep=0.2em}
\setlist[enumerate]{itemsep=0.2em,topsep=0.2em}
\renewcommand{\arraystretch}{1.25}

\newcommand{\Bel}{\operatorname{Bel}}
\newcommand{\pBel}{\overline{\operatorname{Bel}}}
\newcommand{\LM}{\mathrm{L}}

\title{Lösung zu Aufgabe B6.1 (Markow-Lokalisierung)}
\author{}
\date{}

\begin{document}
\maketitle

\section*{Aufgabenstellung und Modell}

Der Roboter bewegt sich deterministisch in einem ringförmigen Korridor mit $16$ Gitterzellen
($x=1,\dots,16$, im Uhrzeigersinn aufsteigend nummeriert). In den Zellen
\[
  \{1,\,4,\,6,\,9,\,13\}
\]
sind Landmarken angebracht. Befindet sich der Roboter in einer Zelle mit Landmarke, so detektiert
er diese mit Wahrscheinlichkeit $70\%$. In allen anderen Zellen meldet er aufgrund der
Sensorungenauigkeit dennoch mit $25\%$ Wahrscheinlichkeit eine Landmarke.

\begin{center}
\setlength{\tabcolsep}{8pt}
\begin{tabular}{|c|c|c|c|c|}
\hline
$1^{\bullet}$ & $2$ & $3$ & $4^{\bullet}$ & $5$ \\\hline
$16$ & & & & $6^{\bullet}$ \\\hline
$15$ & & & & $7$ \\\hline
$14$ & & & & $8$ \\\hline
$13^{\bullet}$ & $12$ & $11$ & $10$ & $9^{\bullet}$ \\\hline
\end{tabular}
\end{center}
{\small ($^{\bullet}$ = Zelle mit Landmarke)}

Der Roboter führt die folgenden $5$ Aktionen aus:
\begin{enumerate}
  \item Der Roboter detektiert eine Landmarke.
  \item Der Roboter bewegt sich $2$ Zellen im Uhrzeigersinn.
  \item Der Roboter detektiert wieder eine Landmarke.
  \item Der Roboter bewegt sich $4$ Zellen im Uhrzeigersinn.
  \item Der Roboter detektiert keine Landmarke.
\end{enumerate}

Gegeben sind damit die folgenden Mess\-wahr\-schein\-lich\-keiten:

\begin{center}
\setlength{\tabcolsep}{10pt}
\begin{tabular}{c|cc}
\toprule
 & $x=\LM$ & $x=\neg\LM$ \\
\midrule
$z=\LM$       & $0{,}7$ & $0{,}25$ \\
$z=\neg\LM$   & $0{,}3$ & $0{,}75$ \\
\bottomrule
\end{tabular}
\end{center}

Für die rekursive Berechnung der Aufenthaltswahrscheinlichkeit verwenden wir die
Prädiktion (a-priori-Schätzung) und die Korrektur (a-posteriori-Schätzung):
\[
  \pBel(x_{t+1}) = \sum_{x_t} P(x_{t+1}\mid x_t,u_t)\cdot \Bel(x_t),
  \qquad
  \Bel(x_{t+1}) = \eta\cdot P(z_t\mid x_{t+1})\cdot \pBel(x_{t+1}),
\]
mit dem Normierungsfaktor
\[
  \eta^{-1} = \sum_{x_{t+1}} P(z_t\mid x_{t+1})\cdot \pBel(x_{t+1}).
\]
Allgemein gilt zu jedem Zeitpunkt
\[
  \Bel(x=1)+\Bel(x=2)+\dots+\Bel(x=16) = 1.
\]
Als Startverteilung nehmen wir die Gleichverteilung an, also
$\Bel_0(x=1)=\dots=\Bel_0(x=16)=\tfrac{1}{16}$.

\section*{Berechnung der Aufenthaltswahrscheinlichkeit (mit vorgegebener Richtung)}

\textbf{1.\ Der Roboter detektiert eine Landmarke.}

Reine Korrektur über die Gleichverteilung. Die Landmarkenzellen werden mit $0{,}7$, alle übrigen
mit $0{,}25$ gewichtet; mit $\eta^{-1}=\tfrac{25}{64}$ ergibt sich
\[
  \Bel_1(x) =
  \begin{cases}
    \tfrac{14}{125}\approx 0{,}112 & x\in\{1,4,6,9,13\},\\[0.3em]
    \tfrac{1}{25}=0{,}04 & \text{sonst.}
  \end{cases}
\]

\textbf{2.\ Der Roboter bewegt sich $2$ Zellen im Uhrzeigersinn.}

Reine Prädiktion. Da die Bewegung deterministisch ist, werden die Wahrscheinlichkeiten lediglich
um $2$ Zellen verschoben:
\[
  \Bel_2(x) =
  \begin{cases}
    0{,}112 & x\in\{3,6,8,11,15\},\\
    0{,}04  & \text{sonst.}
  \end{cases}
\]

\textbf{3.\ Der Roboter detektiert wieder eine Landmarke.}

Korrektur wie in Schritt~1, mit $\eta^{-1}=\tfrac{931}{2500}$. Hier müssen explizit mehrere Fälle
unterschieden werden, je nachdem ob eine verschobene Wahrscheinlichkeit auf eine Landmarkenzelle
trifft:
\[
  \Bel_3(x) =
  \begin{cases}
    \tfrac{4}{19}\approx 0{,}2105 & x=6,\\[0.3em]
    \tfrac{10}{133}\approx 0{,}0752 & x\in\{1,3,4,8,9,11,13,15\},\\[0.3em]
    \tfrac{25}{931}\approx 0{,}0269 & x\in\{2,5,7,10,12,14,16\}.
  \end{cases}
\]

\textbf{4.\ Der Roboter bewegt sich $4$ Zellen im Uhrzeigersinn.}

Wie in Schritt~2 werden die Wahrscheinlichkeiten nur verschoben (um $4$ Zellen):
\[
  \Bel_4(x) =
  \begin{cases}
    0{,}2105 & x=10,\\
    0{,}0752 & x\in\{1,3,5,7,8,12,13,15\},\\
    0{,}0269 & x\in\{2,4,6,9,11,14,16\}.
  \end{cases}
\]

\textbf{5.\ Der Roboter detektiert keine Landmarke.}

Dies verläuft analog zu Schritt~3, lediglich nun mit den Gegenwahrscheinlichkeiten
($0{,}3$ bzw.\ $0{,}75$) und $\eta^{-1}=\tfrac{1203}{1862}$:
\[
  \Bel_5(x) =
  \begin{cases}
    \tfrac{98}{401}\approx 0{,}2444 & x=10,\\[0.3em]
    \tfrac{35}{401}\approx 0{,}0873 & x\in\{3,5,7,8,12,15\},\\[0.3em]
    \tfrac{14}{401}\approx 0{,}0349 & x\in\{1,13\},\\[0.3em]
    \tfrac{25}{802}\approx 0{,}0312 & x\in\{2,11,14,16\},\\[0.3em]
    \tfrac{5}{401}\approx 0{,}0125 & x\in\{4,6,9\}.
  \end{cases}
\]

Dies bedeutet, dass mit der höchsten Wahrscheinlichkeit schlussendlich Zelle~$10$ erreicht wird
($\Bel_5(x=10)=0{,}244$). Rechnet man dies mittels der Aktionen (insgesamt $6$ Zellen im
Uhrzeigersinn) zurück, so korrespondiert dies zu \textbf{Zelle~4 als wahrscheinlichstem
Startpunkt}.

\section*{Berechnung der Aufenthaltswahrscheinlichkeit (mit unbekannter Richtung)}

Da nun keine Angabe über die Bewegungsrichtung vorliegt, gehen wir von einer
$50/50$-Verteilung für die Richtung (im bzw.\ entgegen dem Uhrzeigersinn) aus. Die Korrektur\-schritte
($1$, $3$, $5$) bleiben unverändert; lediglich die Prädiktion wird zu
\[
  \pBel(x_{t+1}) = \tfrac{1}{2}\,\Bel(x_{t+1}-d) + \tfrac{1}{2}\,\Bel(x_{t+1}+d),
\]
da jede Zelle nun aus zwei Vorzuständen erreichbar ist ($d$ = zurückgelegte Zellenzahl).

\textbf{1.\ Der Roboter detektiert eine Landmarke.}

Die Ergebnisse lassen sich $1{:}1$ aus dem ersten Aufgabenteil übertragen:
\[
  \Bel_1(x) =
  \begin{cases}
    \tfrac{14}{125}\approx 0{,}112 & x\in\{1,4,6,9,13\},\\[0.3em]
    \tfrac{1}{25}=0{,}04 & \text{sonst.}
  \end{cases}
\]

\textbf{2.\ Der Roboter bewegt sich $2$ Zellen.}

Grundsätzlich das gleiche Vorgehen, jedoch kann jede Zelle stets aus zwei Vorzuständen erreicht
werden (z.\,B.\ Zelle~$1$ aus Zelle~$3$ und Zelle~$15$):
\[
  \Bel_2(x) =
  \begin{cases}
    0{,}112 & x\in\{11,15\},\\
    \tfrac{19}{250}=0{,}076 & x\in\{2,3,4,6,7,8\},\\
    0{,}04  & x\in\{1,5,9,10,12,13,14,16\}.
  \end{cases}
\]

\textbf{3.\ Der Roboter detektiert wieder eine Landmarke.}

Korrektur mit $\eta^{-1}=\tfrac{931}{2500}$:
\[
  \Bel_3(x) =
  \begin{cases}
    \tfrac{1}{7}\approx 0{,}1429 & x\in\{4,6\},\\[0.3em]
    \tfrac{10}{133}\approx 0{,}0752 & x\in\{1,9,11,13,15\},\\[0.3em]
    \tfrac{5}{98}\approx 0{,}0510 & x\in\{2,3,7,8\},\\[0.3em]
    \tfrac{25}{931}\approx 0{,}0269 & x\in\{5,10,12,14,16\}.
  \end{cases}
\]

\textbf{4.\ Der Roboter bewegt sich $4$ Zellen.}

Ähnlich zu Schritt~2 gibt es je $2$ Vorzustände, aus denen ein Feld erreicht werden kann
(z.\,B.\ Feld~$5$ aus $1$ und $9$):
\[
  \Bel_4(x) =
  \begin{cases}
    \tfrac{79}{931}\approx 0{,}0849 & x\in\{2,8,10,16\},\\[0.3em]
    \tfrac{10}{133}\approx 0{,}0752 & x\in\{5,13\},\\[0.3em]
    \tfrac{235}{3724}\approx 0{,}0631 & x\in\{3,7,11,15\},\\[0.3em]
    \tfrac{5}{98}\approx 0{,}0510 & x\in\{1,9\},\\[0.3em]
    \tfrac{145}{3724}\approx 0{,}0389 & x\in\{4,6,12,14\}.
  \end{cases}
\]

\textbf{5.\ Der Roboter detektiert keine Landmarke.}

Genauso wie Schritt~3, nur mit den Gegenwahrscheinlichkeiten und $\eta^{-1}=\tfrac{249}{392}$:
\[
  \Bel_5(x) =
  \begin{cases}
    \tfrac{158}{1577}\approx 0{,}1002 & x\in\{2,8,10,16\},\\[0.3em]
    \tfrac{140}{1577}\approx 0{,}0888 & x=5,\\[0.3em]
    \tfrac{235}{3154}\approx 0{,}0745 & x\in\{3,7,11,15\},\\[0.3em]
    \tfrac{145}{3154}\approx 0{,}0460 & x\in\{12,14\},\\[0.3em]
    \tfrac{56}{1577}\approx 0{,}0355 & x=13,\\[0.3em]
    \tfrac{2}{83}\approx 0{,}0241 & x\in\{1,9\},\\[0.3em]
    \tfrac{29}{1577}\approx 0{,}0184 & x\in\{4,6\}.
  \end{cases}
\]

Hier sind am wahrscheinlichsten die Zellen $2$, $8$, $10$ oder $16$ als Endposition
(mit $\Bel_5\approx 0{,}1$). Jedoch ist alles allgemein sehr unklar geworden, nicht so eindeutig
wie bei der anderen Aufgabe. Eine Rückrechnung auf den wahrscheinlichsten Initialzustand wäre zwar
grundsätzlich möglich, würde allerdings aufgrund der gewichteten Kanten mit der zusätzlichen
Komplexität den Rahmen etwas sprengen.

\end{document}
